\documentclass[10pt,a4paper]{amsart}
\usepackage{amsmath,amssymb,amsthm,graphicx}
\usepackage{verbatim}
\newenvironment{reference}{\comment}{\endcomment}
%\newenvironment{reference}{}{}
\newenvironment{slogan}{\comment}{\endcomment}
\newenvironment{history}{\comment}{\endcomment}

% For commutative diagrams we use Xy-pic
\usepackage[all]{xy}

% We use 2cell for 2-commutative diagrams.
\xyoption{2cell}
% Theorem environments.
%
\theoremstyle{plain}
\newtheorem{theorem}[subsection]{Theorem}
\newtheorem{proposition}[subsection]{Proposition}
\newtheorem{lemma}[subsection]{Lemma}

\theoremstyle{definition}
\newtheorem{definition}[subsection]{Definition}
\newtheorem{example}[subsection]{Example}
\newtheorem{exercise}[subsection]{Exercise}
\newtheorem{situation}[subsection]{Situation}

\theoremstyle{remark}
\newtheorem{remark}[subsection]{Remark}
\newtheorem{remarks}[subsection]{Remarks}

\numberwithin{equation}{subsection}

% Macros
%
\def\lim{\mathop{\mathrm{lim}}\nolimits}
\def\colim{\mathop{\mathrm{colim}}\nolimits}
\def\Spec{\mathop{\mathrm{Spec}}}
\def\Hom{\mathop{\mathrm{Hom}}\nolimits}
\def\Ext{\mathop{\mathrm{Ext}}\nolimits}
\def\SheafHom{\mathop{\mathcal{H}\!\mathit{om}}\nolimits}
\def\SheafExt{\mathop{\mathcal{E}\!\mathit{xt}}\nolimits}
\def\Sch{\mathit{Sch}}
\def\Mor{\mathop{\mathrm{Mor}}\nolimits}
\def\Ob{\mathop{\mathrm{Ob}}\nolimits}
\def\Sh{\mathop{\mathit{Sh}}\nolimits}
\def\NL{\mathop{N\!L}\nolimits}
\def\CH{\mathop{\mathrm{CH}}\nolimits}
\def\proetale{{pro\text{-}\acute{e}tale}}
\def\etale{{\acute{e}tale}}
\def\QCoh{\mathit{QCoh}}
\def\Ker{\mathop{\mathrm{Ker}}}
\def\Im{\mathop{\mathrm{Im}}}
\def\Coker{\mathop{\mathrm{Coker}}}
\def\Coim{\mathop{\mathrm{Coim}}}

% Boxtimes
%
\DeclareMathSymbol{\boxtimes}{\mathbin}{AMSa}{"02}

%
% Macros for moduli stacks/spaces
%
\def\QCohstack{\mathcal{QC}\!\mathit{oh}}
\def\Cohstack{\mathcal{C}\!\mathit{oh}}
\def\Spacesstack{\mathcal{S}\!\mathit{paces}}
\def\Quotfunctor{\mathrm{Quot}}
\def\Hilbfunctor{\mathrm{Hilb}}
\def\Curvesstack{\mathcal{C}\!\mathit{urves}}
\def\Polarizedstack{\mathcal{P}\!\mathit{olarized}}
\def\Complexesstack{\mathcal{C}\!\mathit{omplexes}}
% \Pic is the operator that assigns to X its picard group, usage \Pic(X)
% \Picardstack_{X/B} denotes the Picard stack of X over B
% \Picardfunctor_{X/B} denotes the Picard functor of X over B
\def\Pic{\mathop{\mathrm{Pic}}\nolimits}
\def\Picardstack{\mathcal{P}\!\mathit{ic}}
\def\Picardfunctor{\mathrm{Pic}}
\def\Deformationcategory{\mathcal{D}\!\mathit{ef}}
\usepackage{cleveref}
\crefname{originalsection}{original section}{original sections}
\Crefname{originalsection}{Original section}{Original sections}
\makeatletter
\newlabel{algebra-section-smooth}{{00T1}{source}{}{sourcepointer-1}{}}
\newlabel{algebra-section-smooth@cref}{{[originalsection][0][]00T1}{[1][0][]source}{}{}{}}
\makeatother
\crefname{originalequation}{original equation}{original equations}
\Crefname{originalequation}{Original equation}{Original equations}
\makeatletter
\newlabel{lemma-equation-plane-curve}{{0BYB}{source}{}{sourcepointer-2}{}}
\newlabel{lemma-equation-plane-curve@cref}{{[originalequation][0][]0BYB}{[1][0][]source}{}{}{}}
\makeatother
\crefname{originalsection}{original section}{original sections}
\Crefname{originalsection}{Original section}{Original sections}
\makeatletter
\newlabel{morphisms-section-universally-bounded}{{03J3}{source}{}{sourcepointer-3}{}}
\newlabel{morphisms-section-universally-bounded@cref}{{[originalsection][0][]03J3}{[1][0][]source}{}{}{}}
\makeatother
\title{An infinite base field supplies a separable closed point of degree at most d on the smooth locus of a plane curve}
\author{The Stacks Project Authors}
\date{stacks-project:curves.tex:local-snapshot-20261012; GNU FDL 1.2 or later}
\begin{document}
\maketitle
\noindent\textit{Editorial note.} All delivered mathematical text is the author's own, copied from the pinned Stacks Project chapter file curves.tex; no statement or proof step was rewritten, reordered or removed.\par\medskip
\noindent\textit{Editorial note.} The counted proof is the author's complete argument: the chart reduction with the nonzero partial derivative, the etale projection, the fibre-degree bound, and the closing separable-point construction are all retained verbatim.\par\medskip
\noindent\textit{Editorial note.} The label lemma-smooth-plane-curve-point-over-separable names the counted statement, and the proof environment immediately following it in the source is delivered verbatim as its proof.\par\medskip
\noindent\textit{Editorial note.} The counted statement is carried by the author's own lemma environment in the chapter source.\par\medskip
\noindent\textit{Allowed standard prerequisites (editorial).}\begin{itemize}
\item The plane-curve setup and the affine chart. Same chapter, Lemma lemma-equation-plane-curve; Algebra, Section algebra-section-smooth. Identifies the chart D\_+(T\_0) with the spectrum of k[x, y]/(f) and, via the Jacobian discussion, makes one partial derivative nonzero at the given smooth point.
\item Etale projections and universally bounded fibres. Morphisms, Section morphisms-section-universally-bounded. Gives the etale projection V -> Spec(k[x]) after shrinking, the openness of its image, and the bound of degree at most d on the fibres since f has degree at most d in y.
\end{itemize}\medskip
\ifdefined\hypertarget\hypertarget{sourcepointer-1}{}\fi
\noindent\textit{Source reference.} Original section 00T1: sources/stacks-chow/tags.txt, source lines 1059--1059. Prerequisite pointer invoked at proof line(s) 1967; it names a standard fact of a sibling chapter that the counted argument uses, and every step of the counted proof built on it is retained verbatim in the delivered body.\par\medskip
\ifdefined\hypertarget\hypertarget{sourcepointer-2}{}\fi
\noindent\textit{Source reference.} Original equation 0BYB: sources/stacks-curves/pinned-original-source.tex, source lines 1859--1859. Prerequisite pointer invoked at proof line(s) 1951; it names a standard fact of the same chapter that the counted argument uses, and every step of the counted proof built on it is retained verbatim in the delivered body.\par\medskip
\ifdefined\hypertarget\hypertarget{sourcepointer-3}{}\fi
\noindent\textit{Source reference.} Original section 03J3: sources/stacks-chow/tags.txt, source lines 4625--4625. Prerequisite pointer invoked at proof line(s) 1978; it names a standard fact of a sibling chapter that the counted argument uses, and every step of the counted proof built on it is retained verbatim in the delivered body.\par\medskip
\begin{lemma}
\label{lemma-smooth-plane-curve-point-over-separable}
Let $Z \subset \mathbf{P}^2_k$ be as in Lemma \ref{lemma-equation-plane-curve}
and let $I(Z) = (F)$ for some $F \in k[T_0, T_1, T_2]$.
If $Z \to \Spec(k)$ is smooth in at least one point and $k$ is infinite, then
there exists a closed point $z \in Z$ contained in the smooth
locus such that $\kappa(z)/k$ is finite separable of degree
at most $d$.
\end{lemma}

\begin{proof}
Suppose that $z' \in Z$ is a point where $Z \to \Spec(k)$ is smooth.
After renumbering the coordinates if necessary we may assume
$z'$ is contained in $D_+(T_0)$. Set $f = F(1, x, y) \in k[x, y]$.
Then $Z \cap D_+(X_0)$ is isomorphic to the spectrum of $k[x, y]/(f)$.
Let $f_x, f_y$ be the partial derivatives of $f$ with respect to
$x, y$. Since $z'$ is a smooth point of $Z/k$ we see that either
$f_x$ or $f_y$ is nonzero in $z'$ (see discussion in
Algebra, Section \ref{algebra-section-smooth}).
After renumbering the coordinates
we may assume $f_y$ is not zero at $z'$. Hence there is a nonempty
open subscheme $V \subset Z \cap D_{+}(X_0)$ such that the
projection
$$
p : V \longrightarrow \Spec(k[x])
$$
is \'etale. Because the degree of $f$ as a polynomial in $y$
is at most $d$, we see that the degrees of the fibres of the
projection $p$ are at most $d$ (see discussion in
Morphisms, Section \ref{morphisms-section-universally-bounded}).
Moreover, as $p$ is \'etale
the image of $p$ is an open $U \subset \Spec(k[x])$.
Finally, since $k$ is infinite, the set of $k$-rational points
$U(k)$ of $U$ is infinite, in particular not empty. Pick any
$t \in U(k)$ and let $z \in V$ be a point mapping to $t$.
Then $z$ works.
\end{proof}

\end{document}
